% !TEX root = main.tex
\chapter*{Sigles i notació\label{Acronyms}}
\addcontentsline{toc}{chapter}{Sigles i notació}
\fancyhead[RE]{\slshape \textsf{Sigles i notació}}
\fancyhead[LO]{\slshape \textsf{Sigles i notació}}

\begin{longtable}{>{\bfseries}p{2.7cm}p{11.8cm}}
ARM & Arquitectura de processadors \textit{Advanced RISC Machines}. \\
bps & Bits per mostra espectral; en aquesta memòria, $\bps$, equivalent a bits
per píxel i banda. \\
CCSDS & \textit{Consultative Committee for Space Data Systems}. \\
CSMR & Implementació de referència del mètode neuronal de compressió
multiràtio i qualitat fixa emprada per generar fitxers \code{.tfci}. \\
DWT & \textit{Discrete Wavelet Transform}. \\
GDN/IGDN & \textit{Generalized Divisive Normalization} i transformada inversa. \\
MSE & \textit{Mean Squared Error}, error quadràtic mitjà. \\
NDVI & \textit{Normalized Difference Vegetation Index}. \\
NDWI & \textit{Normalized Difference Water Index}. \\
OBC & \textit{On-Board Computer}. \\
OpenMP & Interfície de programació per paral·lelisme de memòria compartida. \\
PSNR & \textit{Peak Signal-to-Noise Ratio}. \\
Q-map & Matriu de valors $Q$ que assigna una qualitat a cada bloc espacial. \\
ROI & \textit{Region of Interest}, regió d'interès. \\
RSS & \textit{Resident Set Size}, memòria resident màxima del procés. \\
SWIR & Infraroig d'ona curta. \\
TFC & \textit{TensorFlow Compression}. \\
VAE & \textit{Variational Autoencoder}. \\
\end{longtable}
